% Figure 11.5: two-loop muon magnetic-moment diagram.
% Source: physical PDF p. 512 (printed p. 489).
\begingroup
\renewcommand{\thefigure}{11.\arabic{figure}}
\setcounter{figure}{4}
\begin{figure}[tbp]
\centering
\begin{tikzpicture}[
  electron/.style={line width=0.65pt,
    postaction={decorate},decoration={markings,
      mark=at position 0.10 with {\arrow{Stealth[length=4.5pt,width=3.4pt]}}}},
  muon/.style={line width=2.2pt,
    postaction={decorate},decoration={markings,
      mark=at position 0.28 with {\arrow{Stealth[length=5pt,width=4pt]}}}},
  photon/.style={line width=0.6pt,decorate,
    decoration={snake,amplitude=1.0pt,segment length=5pt}}
]
  \draw[muon] (0,0) -- (0,4.65);
  \draw[photon] (-2.05,2.25) -- (0,2.25);
  \draw[photon] (0,3.75) .. controls (0.60,3.55) and (0.95,2.85) .. (1.45,2.47);
  \draw[photon] (0,0.75) .. controls (0.60,0.90) and (0.95,1.28) .. (1.45,1.35);
  \draw[electron] (1.45,1.35) arc[start angle=210,end angle=-150,x radius=0.75cm,y radius=1.12cm];
\end{tikzpicture}
\caption{A two-loop diagram for the muon magnetic moment. Here the heavy straight line represents a muon; the light wavy lines are photons; and the other light lines are electrons. This diagram makes a relatively large contribution to the fourth-order muon gyromagnetic ratio, proportional to $\ln(m_\mu/m_e)$.}
\label{fig:11.5}
\end{figure}
\endgroup
